% !TEX TS-program = pdflatex
\documentclass[aspectratio=169]{beamer}
\usetheme{Madrid}
\usecolortheme{default}
\usepackage{booktabs}
\setbeamertemplate{navigation symbols}{}
\setbeamertemplate{footline}[frame number]

\title{Free the tail, keep the head}
\subtitle{Tail-only preemption in vLLM $\cdot$ Qwen3-14B $\cdot$ one H200 $\cdot$ vLLM 0.11.1}
\author{Sagnik Biswas}
\date{22 October 2026}

\begin{document}

\begin{frame}
\titlepage
\end{frame}

\begin{frame}{The bug: 7 minutes becomes 31}

Pool at 45\% memory: \textbf{29.5 GiB} (12{,}100 blocks of 16 tokens).
Each request: 8{,}064-token prompt, forced to 32{,}768 output tokens. Caching off.

\begin{center}
\begin{tabular}{lrrr}
\toprule
run & KV needed & preemptions & time \\
\midrule
4 requests & 24.9 GiB & 0 & 7.1 min \\
5 requests & 31.2 GiB & 807 & 30.7 min \\
\bottomrule
\end{tabular}
\end{center}

The fifth request overflows the pool by 1.7 GiB, about 6\%.
Stock vLLM answers by throwing away a \emph{whole} request and recomputing it, over and over.

\end{frame}

\begin{frame}{What stock vLLM does on preemption}

\begin{itemize}
  \item Free \emph{every} KV block of the victim.
  \item Set its progress to zero and put it back in the queue.
  \item On resume, prefill the prompt and everything it has generated so far.
\end{itemize}

\vspace{0.5em}
The pool was short by a few hundred blocks. The victim held thousands.
Most of what gets freed did not need to be freed.

\end{frame}

\begin{frame}{The patch: a scheduler subclass}

\texttt{serving/tail\_preempt\_scheduler.py}, passed in through \texttt{scheduler\_cls}. No vLLM fork.

\begin{itemize}
  \item On preemption, keep the victim's head blocks pinned. Free nothing yet.
  \item When someone needs memory, trim the pinned blocks one at a time from the \emph{tail} of the lowest-priority victim.
  \item On resume, hand the surviving head back. Recompute only the trimmed tail.
  \item Wait to resume until the request's remaining blocks fit, so it is not preempted again right away.
  \item Fall back to stock behavior when something is unsupported (spec decode, KV connectors, multiple KV groups, async scheduling).
\end{itemize}

Works with prefix caching on or off.

\end{frame}

\begin{frame}{Five requests, same pod}

\begin{center}
\begin{tabular}{llrrr}
\toprule
caching & scheduler & preemptions & recomputed & time \\
\midrule
off & stock & 815 & 31.5M & 30.2 min \\
off & \textbf{tail} & \textbf{1} & \textbf{8.4k} & \textbf{8.0 min} \\
on & stock & 1 & 8.4k & 8.0 min \\
on & tail & 1 & 8.4k & 8.0 min \\
\midrule
\multicolumn{2}{l}{earlier: 4 requests, no pressure} & 0 & 0 & 7.1 min \\
\bottomrule
\end{tabular}
\end{center}

Every arm reached all 32{,}768 tokens. With caching off, tail is \textbf{3.8$\times$ faster} and recomputes 3{,}700$\times$ fewer tokens.
With caching on, stock and tail are identical: same one preemption, same 528 blocks recomputed.

\end{frame}

\begin{frame}{Inside the tail run}

\begin{center}
\begin{tabular}{lr}
\toprule
event & count \\
\midrule
preemptions & 1 \\
blocks pinned at preemption & 2{,}419 \\
blocks actually trimmed & 528 \\
tokens reused on resume & 30{,}256 \\
tokens recomputed & 8{,}449 \\
\bottomrule
\end{tabular}
\end{center}

The pool needed 528 blocks. Tail freed 528 blocks, not 2{,}419.
Stock freed whole requests 815 times and recomputed 31.5M tokens.

\end{frame}

\begin{frame}{The honest caveat}

Prefix caching is on by default in vLLM, and it already rescues this case:
stock with caching on took \textbf{478 s with 1 preemption}, the same as tail.

\begin{itemize}
  \item Freed blocks keep their hashes, and the free queue evicts the tail first.
    So stock plus caching recovers most of the head on its own.
  \item The 31-minute behavior needs caching off.
  \item What tail adds: the same recovery \emph{without} relying on prefix caching.
    That matters where caching is off or does not apply (prompt embeddings on vLLM 0.11.0, multimodal inputs, per-request salts).
\end{itemize}

With caching on, tail adds nothing: stock already recovers exactly the same head.

\end{frame}

\begin{frame}{LatentMAS on AIME: no gain}

60 AIME 2024+2025 problems, up to 20k output tokens, the Judger served through vLLM.

\begin{center}
\begin{tabular}{rlrrr}
\toprule
concurrency & scheduler & output tok/s & preemptions & recompute tokens \\
\midrule
16 & stock & 844 & 0 & 0 \\
16 & tail & 874 & 0 & 0 \\
32 & stock & 1033 & 35 & 94k \\
32 & tail & 999 & 62 & 149k \\
60 & stock & 996 & 54 & 240k \\
60 & tail & 1002 & 51 & 208k \\
\bottomrule
\end{tabular}
\end{center}

Differences are inside the 3.5\% noise floor. Recompute is a cheap prefill here, at most about 4\% of time.

\end{frame}

\begin{frame}{Takeaways}

\begin{itemize}
  \item Under real overflow with caching off, tail preemption turns 815 preemptions and 30.2 min into 1 preemption and 8.0 min.
  \item With caching on, stock gets exactly the same result. Tail only helps where caching is off or does not apply.
  \item On LatentMAS AIME serving, recompute is too cheap for either scheduler to matter.
  \item Next: workloads where recompute is expensive and caching does not help,
    such as long contexts, many victims at once, or caching off by design.
\end{itemize}

\end{frame}

\end{document}
